% ============================================================================
% Fichier  : partie3/P03_C14_auditer_risques.tex
% Titre    : Auditer la Gestion des Risques et le SMSI
% Auteur   : MINKA MI NGUIDJOÏ Thierry Emmanuel
% Version  : 1.0 -- novembre 2026
% ============================================================================
% Sources : ISO 31000:2018 ; ISO/IEC 27001:2022 (clauses 6, 8, 9, 10) ;
% ISO/IEC 27007:2020 ; IIA, modèle des trois lignes (2020). Audit d'octobre
% 2026 de LiZbiZ : cas fictif.
% ============================================================================

\chapter{Auditer la Gestion des Risques et le SMSI}
\label{chap:auditer-risques}
\chaptermark{Auditer les risques et le SMSI}

\epigraph{%
    « Un registre qui n'a jamais été surpris par un incident\\
    n'a probablement jamais été confronté à un incident. »%
}{Principe du cours}

\niveauChapitre{Intermédiaire à expert}{Chapitres~\ref{chap:traiter},
\ref{chap:audit} et~\ref{chap:conduire}.}

\begin{boxobjectifs}
\begin{itemize}
    \item Auditer la conception, le fonctionnement et l'efficacité d'un
          dispositif de gestion des risques.
    \item Rapprocher les incidents du registre pour trouver les risques
          oubliés.
    \item Tester la reproductibilité des cotations et la réalité des
          mesures déclarées.
    \item Auditer un indicateur clé de risque~: définition, population,
          source.
    \item Relier l'audit aux clauses de l'ISO/IEC~27001.
\end{itemize}
\end{boxobjectifs}

\bigskip

% ============================================================================
\section*{Question directrice : le registre de LiZbiZ dit-il vrai~?}
% ============================================================================

En octobre, Mme~NKENG audite le dispositif mis en place par Mme~MOUAFO
depuis avril. Le registre compte 46~risques, le plan de traitement six
actions, le tableau de bord quatre indicateurs, tous verts. Tout semble en
ordre. L'audit doit dire si le dispositif est bien conçu, s'il fonctionne,
et s'il produit l'effet attendu.

\medskip
\textbf{Plan du chapitre.}
\begin{enumerate}
    \item Trois niveaux d'examen (section~\ref{sec:c14-niveaux}).
    \item Rapprocher incidents et registre (section~\ref{sec:c14-incidents}).
    \item Tester cotations et mesures (section~\ref{sec:c14-tests}).
    \item Auditer un indicateur (section~\ref{sec:c14-indicateur}).
    \item Les clauses de l'ISO/IEC 27001 (section~\ref{sec:c14-27001}).
\end{enumerate}

\separateur

% ============================================================================
\section{Trois Niveaux d'Examen}
\label{sec:c14-niveaux}
% ============================================================================

\begin{tableaucours}{L{2.8cm}L{4.4cm}L{5.0cm}}{%
    \tableauHeader{Niveau} &
    \tableauHeader{Question} &
    \tableauHeader{Preuves}
}
Conception & Le dispositif est-il bien pensé~? & Politique, appétence,
critères, rôles, méthode \\
Fonctionnement & Est-il appliqué comme prévu~? & Registre à jour, comptes
rendus du comité, plans suivis, indicateurs mesurés \\
Efficacité & Atteint-il son but~? & Incidents prévus ou non, appétence
respectée, risques réellement réduits \\
\end{tableaucours}
\vspace{-0.8em}
\begin{center}{\small\itshape Tableau~14.1. Auditer un dispositif de risques}\end{center}

\begin{boxalert}
\textbf{Violation typique.} S'arrêter à la conception (« la méthode existe,
elle est conforme à l'ISO~31000 »).

\smallskip
\textbf{Règle~:} une méthode excellente, mal appliquée, ne protège de rien.
L'auditeur descend jusqu'à l'efficacité~: les risques ont-ils baissé, et
le registre a-t-il prévu les incidents survenus~?
\end{boxalert}

% ============================================================================
\section{Rapprocher Incidents et Registre}
\label{sec:c14-incidents}
% ============================================================================

Chaque incident survenu est un test du registre~: figurait-il parmi les
scénarios identifiés~? Avec quelle cote~? Un incident absent du registre
révèle un trou d'identification~; un incident d'un risque coté
« acceptable » peut révéler une sous-estimation.

\begin{tableaucours}{L{4.8cm}L{1.6cm}L{5.8cm}}{%
    \tableauHeader{Incident du trimestre} &
    \tableauHeader{Au registre} &
    \tableauHeader{Lecture}
}
Coupure du lien de Garoua (19~h) & R-03 & Prévu, cote cohérente \\
Hameçonnage d'une assistante, signalé à temps & R-02 & Prévu \\
Erreur de tarif sur un contrat & R-19 & Prévu \\
Vol d'un ordinateur portable d'un commercial & Non & Trou~: pas de
risque « perte d'équipement nomade » \\
Panne de la climatisation du centre de Douala & Non & Trou~: dépendances
techniques des centres non examinées \\
\end{tableaucours}
\vspace{-0.8em}
\begin{center}{\small\itshape Tableau~14.2. Rapprochement des incidents de juillet à septembre}\end{center}

% ============================================================================
\section{Tester Cotations et Mesures}
\label{sec:c14-tests}
% ============================================================================

\begin{boxprocedure}
\textbf{Quatre tests du dispositif}
\begin{enumerate}
    \item \textbf{Reproductibilité}~: faire recoter cinq risques par deux
          personnes, séparément~; comparer.
    \item \textbf{Réalité des mesures}~: pour chaque action déclarée
          réalisée, demander la preuve de mise en œuvre et de vérification.
    \item \textbf{Acceptations}~: vérifier que chaque risque résiduel
          accepté l'a été par écrit, par la bonne personne.
    \item \textbf{Propriétaires}~: interroger trois propriétaires sur leur
          risque, son plan et son indicateur.
\end{enumerate}
\end{boxprocedure}

\begin{boxscenario}{Opération LIZBIZ. Ce que l'audit d'octobre a trouvé}
\textbf{Constatations.}
\begin{itemize}
    \item Recotation~: trois risques sur cinq diffèrent d'au moins un
          niveau entre les deux cotateurs.
    \item Plan de traitement~: six actions~; quatre réalisées et
          vérifiées~; une en retard (second lien de Garoua)~; une déclarée
          réalisée mais contredite par un essai (second facteur, compte
          « vpn\_secours », Chapitre~\ref{chap:conduire}).
    \item Acceptations~: deux risques résiduels acceptés oralement, sans
          trace.
    \item Propriétaires~: deux sur trois connaissent leur plan~; un ignore
          qu'il est propriétaire.
\end{itemize}

\medskip
\textbf{Analyse} (confiance élevée). Le dispositif est bien conçu et
fonctionne en partie~; son efficacité est surestimée par un registre qui
compte comme réalisée une action incomplète.
\end{boxscenario}

% ============================================================================
\section{Auditer un Indicateur}
\label{sec:c14-indicateur}
% ============================================================================

Un indicateur vert peut être faux de trois manières~: la \emph{définition}
ne mesure pas le risque, la \emph{population} est incomplète, la
\emph{source} n'est pas fiable.

\begin{boxobjet}{Constat A-05 · Indicateur du second facteur}
\textbf{Critère~:} indicateur clé de risque « part des connexions
distantes avec second facteur » (Tableau~11.4). \textbf{Constat~:}
l'indicateur affiche 100~\%, mais il est calculé sur les comptes
\emph{nominatifs} seulement~; le compte « vpn\_secours » en est exclu par
construction. Sur toutes les connexions de septembre, le taux réel est de
$(1\,340 - 46) / 1\,340 \approx 96{,}6~\%$, soit l'orange.
\textbf{Preuve~:} requête de calcul de l'indicateur, export du journal.
\textbf{Cause~:} population de l'indicateur définie par la même liste que
la politique, exceptions comprises. \textbf{Classement~:} non-conformité
mineure (l'écart principal est traité au constat A-03).
\end{boxobjet}

\begin{boxmetier}{Auditeur des systèmes d'information}
\textbf{Recalculer un indicateur à partir de la source.} L'auditeur ne
recopie pas le tableau de bord~: il refait le calcul d'au moins un
indicateur depuis les données brutes, et compare.
\end{boxmetier}

% ============================================================================
\section{Les Clauses de l'ISO/IEC 27001}
\label{sec:c14-27001}
% ============================================================================

Pour un système de management de la sécurité, l'audit du dispositif de
risques s'appuie sur les clauses de l'ISO/IEC~27001~\cite{iso27001_2022}
et sur les recommandations de l'ISO/IEC~27007~\cite{iso27007_2020}.

\begin{tableaucours}{L{2.0cm}L{4.4cm}L{5.8cm}}{%
    \tableauHeader{Clause} &
    \tableauHeader{Exigence} &
    \tableauHeader{Preuve demandée chez LiZbiZ}
}
6.1.2 & Appréciation des risques & Méthode, critères, registre daté \\
6.1.3 & Traitement des risques & Plan de traitement, déclaration
d'applicabilité, acceptation des résiduels \\
8.2 et 8.3 & Appréciation et traitement réalisés à intervalles planifiés
& Comptes rendus des revues trimestrielles \\
9.1 & Surveillance et mesure & Définition et calcul des indicateurs \\
9.2 & Audit interne & Programme et rapports d'audit \\
9.3 & Revue de direction & Compte rendu de la revue \\
10.2 & Non-conformité et action corrective & Suivi des constats A-03 et
A-05 \\
\end{tableaucours}
\vspace{-0.8em}
\begin{center}{\small\itshape Tableau~14.3. Clauses de l'ISO/IEC 27001 relatives aux risques}\end{center}

\begin{boxmetier}{Gestionnaire de risques}
\textbf{Accueillir l'audit comme une aide.} Mme~MOUAFO fournit les sources
brutes, pas seulement les tableaux~; un constat sur son indicateur
améliore son dispositif avant qu'un incident ne le fasse.
\end{boxmetier}

\separateur

\begin{boxresume}
\textbf{Ce qu'il faut retenir de ce chapitre~:}
\begin{itemize}
    \item \textbf{Trois niveaux}~: conception, fonctionnement, efficacité.
    \item \textbf{Chaque incident} teste le registre~: prévu ou non, et
          avec quelle cote.
    \item \textbf{Quatre tests}~: reproductibilité, réalité des mesures,
          acceptations, propriétaires.
    \item \textbf{Un indicateur} se vérifie par sa définition, sa
          population et sa source.
    \item \textbf{L'ISO/IEC~27001} (6.1, 8, 9, 10) donne les critères pour un
          système de management.
\end{itemize}
\end{boxresume}

\begin{boxinfo}
\textbf{Points de contrôle de l'élève.} Je sais~: distinguer conception,
fonctionnement et efficacité~; rapprocher incidents et registre~; conduire
les quatre tests~; auditer un indicateur~; citer les clauses de
l'ISO/IEC~27001 relatives aux risques.
\end{boxinfo}

\bigskip

\begin{boxexo}[title={\ding{45}\quad Exercice~14.1. Quel niveau~? \niveaubase}]
Classez en conception, fonctionnement ou efficacité~: « la méthode n'a pas
d'échelle d'impact »~; « le comité ne s'est pas réuni au troisième
trimestre »~; « deux incidents n'étaient pas au registre ».
\end{boxexo}

\begin{boxexo}[title={\ding{45}\quad Exercice~14.2. Trous \niveaubase}]
Formulez les deux risques manquants révélés par le Tableau~14.2.
\end{boxexo}

\begin{boxexo}[title={\ding{45}\quad Exercice~14.3. Recalculer \niveaubase}]
Vérifiez le taux réel de connexions avec second facteur du constat A-05.
\end{boxexo}

\begin{boxexo}[title={\ding{45}\quad Exercice~14.4. Reproductibilité \niveaumoyen}]
Trois risques sur cinq sont recotés différemment. Quelles causes
possibles, et quelles corrections~?
\end{boxexo}

\begin{boxexo}[title={\ding{45}\quad Exercice~14.5. Acceptation orale \niveaumoyen}]
Pourquoi une acceptation orale d'un risque résiduel est-elle un écart~?
Quel critère méconnaît-elle~?
\end{boxexo}

\begin{boxexo}[title={\ding{45}\quad Exercice~14.6. Le propriétaire qui s'ignore \niveaumoyen}]
Un directeur ignore qu'il est propriétaire d'un risque. Rédigez le
constat et sa recommandation.
\end{boxexo}

\begin{boxexo}[title={\ding{45}\quad Exercice~14.7. Population \niveaumoyen}]
Donnez un autre exemple d'indicateur vert faussé par sa population, pris
dans la sécurité des postes de travail.
\end{boxexo}

\begin{boxexo}[title={\ding{45}\quad Exercice~14.8. Clause 10.2 \niveaumoyen}]
Que doit contenir, selon la clause 10.2, la réponse de LiZbiZ au constat
A-03~?
\end{boxexo}

\begin{boxexo}[title={\ding{45}\quad Exercice~14.9. Programme de tests \niveauexpert}]
Construisez la liste de vérification complète de l'audit du dispositif de
risques de LiZbiZ~: points, critère, preuve, test.
\end{boxexo}

\begin{boxexo}[title={\ding{45}\quad Exercice~14.10. Opinion d'ensemble \niveauexpert}]
Rédigez la conclusion d'ensemble de Mme~NKENG sur le dispositif, en trois
paragraphes~: conception, fonctionnement, efficacité.
\end{boxexo}

\bigskip

\begin{boxinfo}
\textbf{Pour aller plus loin.}
\begin{itemize}
    \item ISO/IEC~27001~\cite{iso27001_2022}~; ISO/IEC~27007~\cite{iso27007_2020}.
    \item IIA, modèle des trois lignes~\cite{iia2020troislignes}.
\end{itemize}
\end{boxinfo}

\begin{boxinfo}
\textbf{Transition.} Les constats doivent être écrits, présentés et
suivis. Le Chapitre~\ref{chap:rapport} traite du rapport et du suivi des
actions.
\end{boxinfo}
